\documentclass[11pt,reqno]{amsart}
\usepackage[T1]{fontenc}
\usepackage{lmodern,amsmath,amssymb,mathtools,microtype,booktabs,array}
\usepackage[margin=1in]{geometry}
\usepackage[hidelinks]{hyperref}
\hypersetup{pdftitle={General Theta Foundations I: Acquired Geometry and Causal Resource Transfer},pdfauthor={Qian Qi}}
\ifdefined\pdfinfoomitdate\pdfinfoomitdate=1\fi
\ifdefined\pdftrailerid\pdftrailerid{}\fi
\ifdefined\pdfsuppressptexinfo\pdfsuppressptexinfo=15\fi
\numberwithin{equation}{section}
\newtheorem{theorem}{Theorem}[section]
\newtheorem{proposition}[theorem]{Proposition}
\newtheorem{lemma}[theorem]{Lemma}
\newtheorem{corollary}[theorem]{Corollary}
\theoremstyle{definition}
\newtheorem{definition}[theorem]{Definition}
\newtheorem{example}[theorem]{Example}
\theoremstyle{remark}
\newtheorem{remark}[theorem]{Remark}
\newcommand{\E}{\mathbb E}
\newcommand{\Pp}{\mathbb P}
\newcommand{\R}{\mathbb R}
\newcommand{\one}{\mathbf1}
\newcommand{\tr}{\operatorname{tr}}
\newcommand{\law}{\operatorname{Law}}
\newcommand{\TV}{\mathrm{TV}}
\newcommand{\cp}{\mathrm{cp}}
\newcommand{\aut}{\mathrm{aut}}
\newcommand{\eps}{\varepsilon}
\newcommand{\Qm}{\mathcal Q}
\DeclareMathOperator*{\argmin}{arg\,min}
\emergencystretch=2em
\title[Acquired geometry and causal resource transfer]{General Theta Foundations I:\\Acquired Geometry and Causal Resource Transfer}
\author{Qian Qi}
\date{October 2026}
\subjclass[2020]{Primary 60G40, 62C20; Secondary 28A80, 60J05, 81P15, 94A29}
\keywords{causal experiment, predictive quotient, acquired geometry, calibrated risk, stopping, finite memory, state-dependent heralding}
\begin{document}
\begin{abstract}
We prove a resource--risk transfer theorem for prepared causal experiments whose acquisition laws may depend on the entire observed history and on an unknown fixed model. An exact common finite observable driver is not assumed. A predictable effort inequality bounds how much productive acquisition any legal strategy can obtain; a finite-state acquisition program supplies the converse at a separately certified cost. Conditional acquired-mass bounds then transfer the same marked quantization profile to checkpoint and autonomous risk, with all live and terminal states charged. We also characterize the existence of a history-uniform quadratic calibration mark by a comparison between a conditional minimax radius and the least residual variance on each actual task fiber. The inherited marked quantization calculation yields a matching waiting--memory--risk law for history-dependent, state-dependent heralding, including singular supports and a retained atomic transition. Controlled hidden-state kernels and noncommuting quantum instruments verify the same result without a common primary report row. Physical-command compatibility replaces action-name counting. A terminal-measure-preserving realization theorem on standard Borel spaces and an explicit finite-bit stopped certificate complete the causal and numerical layers. Uniformity in rare heralding is stated on its feasible rescaled budget, not at a fixed budget below the cost of one acquisition.
\end{abstract}
\maketitle
\section{The problem and the mechanism}
An experiment does not supply a geometric parameter space together with a free observation map. It supplies a preparation, physical commands, causal reports, and executable future tests. A visible history determines predictions of those tests. The resulting predictive quotient may be infinite, singular, or model dependent. Even when the terminal task is finite dimensional, a favorable task value need not be obtainable at a prescribed acquisition cost or executable with a prescribed finite controller.

The problem considered here is to connect these objects in their physical order:
\[
 \begin{gathered}
 \text{prepared kernel}\longrightarrow\text{future tests}
 \longrightarrow\text{predictive quotient}\\
 \longrightarrow\text{causal acquisition}
 \longrightarrow\text{acquired geometry}\longrightarrow\text{risk}.
 \end{gathered}
\]
Three distinctions are essential. A law obtained by selecting favorable reports is not a free preparation. A task proxy is not necessarily the full predictive state. The labels used to acquire a report are not available a second time as terminal prediction values.

A finite observable projection can make exact stopped occupation tractable, but it is restrictive: two histories with the same protocol phase may induce different next-report laws. We replace this exact closure by two inequalities verified before an observer or a budget is chosen. The first is a predictable lower certificate for the physical effort needed to obtain productive reports. The second is a uniform upper certificate for an explicitly executable acquisition program. Both allow the microscopic state, the predictive coordinate and the report law to retain the entire relevant past. Conditional comparison with one actually acquired reference law supplies the remaining geometric control.

Theorem~\ref{thm:main} is the central result. It bounds the risks of all legal history-dependent procedures below, and constructs one parameter-independent finite-state procedure for the upper bound. Its two profiles have different, explicit budget rescalings. They become matching curves when the acquired geometry has the stated regularity; they are not silently identified at an atomic critical point. Theorem~\ref{thm:herald} applies the theorem to genuinely state-dependent heralding, with no independence of successive attempts and no common finite report row. The unknown heralding rate may change with hidden state and history inside certified bounds. Theorem~\ref{thm:gaussian-herald} also derives the ordinary hidden-sign risk without a supplied reliability coordinate, giving a stretched-logarithmic acquisition law for endogenous waiting.

A second issue concerns the origin of the task geometry. Theorem~\ref{thm:fiber} gives a necessary and sufficient criterion for a \emph{uniform quadratic mark of the declared kind}. On each actual full-information task fiber, form the least residual variance and the minimax quadratic radius. Such a mark exists precisely when the radius is bounded by a fixed multiple of that least variance. This is not a criterion for every possible compression method, but it identifies exactly the obstruction to the conditional comparison used here. It also makes equivalent marks interchangeable up to fixed constants.

The physical controller is handled in Section~\ref{sec:memory}. Commands are identified by their physical interface, including legality, not merely by their printed names. Mandatory incompatible physical command sets give a live-state lower bound. This bound, and the optimal live width itself, are invariant under stateless resource-preserving recodings. A recoding that needs to remember a protocol phase is not stateless and its state is charged. Appendix~\ref{sec:flow} separately extends terminal-submeasure-preserving killed-flow realization from finite drivers to standard Borel state spaces. That realization may require the full predictive state; it is never presented as a free finite-memory implementation.

Conditional-envelope rejection was already used in the causal-stopping and causal-acquisition companions~\cite{GTF46,GTF51}; it is not claimed as a new sampling method here. The marked quantization calculation and conditional projection mechanism are inherited from the stopped-experiment and calibrated-protocol companions~\cite{GTF52,GTF53} and are re-proved where used. Classical occupation-measure Markovization is adjacent to, and not invented by, the present work: finite extreme occupations in Borel absorbing models are studied by Piunovskiy and Zhang~\cite{PZExtreme}; strategic-to-occupation continuity and absorption by the same authors~\cite{PZContinuity}; and stationary or chattering sufficiency by Dufour and Prieto-Rumeau~\cite{DPAbsorb,DPChatter}. Our main risk proof does not rely on compactness of the space of all stopped trajectories. It uses predictable count domination and an explicit finite-state acquisition construction.

Finite observer memory alone does not make general partially observed stopping easy. Optimal-cost almost-sure reachability~\cite{CCGK} and fixed-memory reachability~\cite{ACSS} have quite different computational difficulties. The theorem below uses substantive task, mass and effort certificates; it is not an efficient solution of arbitrary POMDPs. The new point relative to the companions is the combination of raw-history-dependent acquisition, a sharp calibration criterion, physical-command accounting, and fully paid implementation. No claim of an exhaustive priority classification or of a solution to every causal experiment is made.

\section{Raw experiments, task fibers and paid acquisition}
\label{sec:setup}
All spaces are standard Borel. A positive normalized causal kernel, a preparation and a fixed bounded terminal audit $Z\in H$ are given, where $H\subset\R^k$ is compact and convex. The model $\theta\in\Theta$ is fixed along a trajectory and unavailable to the observer. The hidden state may include its microscopic past. The complete visible filtration contains every reported failure, every previously used private coin, and every recorded physical cost. A fresh routing coin is conditionally independent of the future audit. Every legal terminating procedure has a stopping time $1\le\tau<\infty$ almost surely under every admitted model; the infimum is restricted to such procedures.

At raw call $t$ a physical command is selected from the already available information. The nonnegative reported cost $c_t$ satisfies $c_t\ge c_*>0$. A predictable flag $A_t\in\{0,1\}$ says whether this call is productive. Only a productive report permits commitment. Its visible task mark is $X_t\in\mathsf X$; an additional report $V_t$ can be informative and may be used by arbitrary competing procedures. The audit is executed after commitment, with no undeclared intervening transition. Preparation and audit costs not included in $c_t$ are displayed separately when a total budget is used.

Fix a task baseline $B_0\ge0$. After the entire current report let
\[
 \mu=\E[Z\mid\mathcal H],\qquad s=\E[\|Z-\mu\|^2\mid\mathcal H]-B_0.
\]
We assume $s\ge0$ on stop-allowed histories. For each mark $x$, let $\mathcal F_x\subset H\times[0,S]$ be the closure of the attainable pairs $(\mu,s)$ over all admitted models and legal full histories with that mark, excluding zero-evidence histories. It is a task fiber, not a set chosen by the encoder. Conditional versions on their actual exposure laws suffice. For the pointwise fiber criterion below we require nonempty compact fibers and Borel functions
\[
 \underline s(x)=\inf_{(u,v)\in\mathcal F_x}v,
 \qquad g_x(r)=\sup_{(u,v)\in\mathcal F_x}\{v+\|u-r\|^2\}
\]
for each $r$ in a fixed countable dense subset of $H$. In applications an explicitly verified compact invariant family can replace the actual fibers for a \emph{sufficient} certificate; failure on a larger family is not a necessity statement about the actual experiment.

The unique minimizer and its value are
\begin{equation}\label{eq:canonical-mark}
 f_*(x)=\argmin_{r\in H}g_x(r),\qquad w_*(x)=g_x(f_*(x)).
\end{equation}
They are Borel, as proved in Theorem~\ref{thm:fiber}. Its raw calibration criterion is
\begin{equation}\label{eq:fiber-ratio}
 w_*(x)\le K\underline s(x),\qquad 1\le K<\infty.
\end{equation}
It implies the uniform comparison
\begin{equation}\label{eq:task-compare}
 k_-\{w(x)+\|f(x)-r\|^2\}
 \le \E[\|Z-r\|^2\mid\mathcal H]-B_0
 \le k_+\{w(x)+\|f(x)-r\|^2\},
\end{equation}
with $(f,w)=(f_*,w_*)$, $k_-=(3K)^{-1}$ and $k_+=2$. Equivalently one may verify any bounded Borel $f:\mathsf X\to H$, $w\ge0$ and positive $k_-,k_+$ satisfying \eqref{eq:task-compare}. For example
\begin{equation}\label{eq:ordinary-calibration}
 a_vw\le s\le A_vw,\qquad \|\mu-f\|^2\le a_bw
\end{equation}
gives $k_-=[\max\{2,(1+2a_b)/a_v\}]^{-1}$ and $k_+=\max\{2,A_v+2a_b\}$. All comparisons are after $V_t$; hence they survive selection using $V_t$ and an independent fresh stopping coin.

\subsection*{The acquired law and its conditional envelopes}
Fix a known reference law $\nu$ of the mark produced by a declared, actually executable reference preparation and acquisition. Its acquisition is not repeated for free during the online task. A known reference member of a model family may be used; it is not the unknown true member. Assume constants $0<\ell\le1\le u<\infty$ such that, for every model and every surviving productive history and command,
\begin{equation}\label{eq:mass-upper}
 \law_\theta(X_t\mid\mathcal H_{t-1},\text{command})\le u\nu.
\end{equation}
This is a measure inequality on all Borel sets. The actual conditional laws can differ for every history.

A parameter-independent acquisition program $\mathcal G$ has $B$ live labels, a designated return label, and compatible physical commands. From that return label it runs until a productive report; after rejection it returns to that label without resetting the hidden physical state. It stores no other history, seed, real posterior or clock. At its productive calls require also
\begin{equation}\label{eq:mass-lower}
 \law_\theta(X_t\mid\mathcal H_{t-1},\text{command})\ge\ell\nu.
\end{equation}
The current mark and the chosen terminal cell can be evaluated from its live label and the current report. Every attempt completes almost surely and has conditional expected raw cost at most $\kappa_+$, uniformly over its starting history and model. These are verified properties of an actual finite-state program, not properties assigned to an abstract state graph.

\subsection*{A lower effort certificate}
There is a bounded Borel potential $V(h)$ on raw continued histories, independent of the true model, equal to zero initially and immediately after every productive call. For every legal pre-call history and command,
\begin{equation}\label{eq:effort-drift}
 \E_\theta[c_t+V(\mathcal H_t)-V(\mathcal H_{t-1})
            \mid\mathcal H_{t-1},\text{command}]
 \ge\kappa_- A_t,\qquad \kappa_->0.
\end{equation}
$V$ may be negative. At commitment its value is zero. A limiting version with uniformly integrable stopped potentials is possible, but boundedness is the hypothesis used here. This certificate constrains \emph{all} admissible strategies, not just $\mathcal G$.

\begin{lemma}[Verification of an acquisition program]\label{lem:navigator}
The upper attempt-cost property follows if a nonnegative function $W(s,h)$ on the product of its live labels and raw predictive histories satisfies $W(s_0,h)\le\kappa_+$ at each return history and
\[
 \E_\theta[c_t+W(S_t,\mathcal H_t)\one_{\{\text{attempt continues}\}}
      \mid S_{t-1},\mathcal H_{t-1}]
 \le W(S_{t-1},\mathcal H_{t-1})
\]
for every surviving attempt history. The proof potential is not stored by the machine.
\end{lemma}
\begin{proof}
Stop at the earlier of attempt completion and $n$ further calls. Conditional telescoping, with the nonnegative remaining potential retained, bounds expected cost by the initial $W$. Monotone convergence bounds the full sum. Since each call costs at least $c_*$, the expected number of calls is finite and completion is almost sure. This argument is uniform over the starting history. It assumes neither regeneration nor a positive completion hazard on every raw call.
\end{proof}

Let $R_\cp^\Theta(N,m)$ be the infimum of worst-model risk over arbitrary legal history-dependent procedures with a fixed dictionary of at most $m$ terminal values and $\sup_\theta\E_\theta\sum_{t\le\tau}c_t\le N$. Let $R_\aut^\Theta(N,M)$ impose a total alphabet of at most $M$ disjoint live and terminal labels. Actions depend only on the live label; routing uses that label and the current command and report. The routing class at this level consists of Borel stochastic kernels. This is an alphabet-constrained model, not a claim that every Borel rule has finite program description or finite arithmetic cost. Those resources require a computable cell certificate as in Proposition~\ref{prop:finite-cells}, or a direct finite-bit construction as in Section~\ref{sec:numerics}. Infeasible risks are $+\infty$. The same $B_0$ and the same audit occur in both risks.

\section{The predictable acquired-geometry theorem}
\label{sec:main}
For a dictionary $C=(r_1,\ldots,r_m)\in H^m$ put
\[
 d_C(x)=w(x)+\min_j\|f(x)-r_j\|^2.
\]
The acquired marked profile is
\begin{equation}\label{eq:profile}
 \Qm_m(b)=\min\left\{\int d_C h\,d\nu:
       C\in H^m,\ 0\le h\le b,\ \int h\,d\nu=1\right\},\qquad b\ge1.
\end{equation}
It is $+\infty$ for $b<1$. Here $h\nu$ is a candidate selected \emph{probability measure}, not a free preparation. Its actual acquisition cost enters through $b$. The minimum is over Borel representatives modulo $\nu$-null sets.

\begin{theorem}[Predictable acquired geometry with paid causal acquisition]\label{thm:main}
Consider the raw experiment of Section~\ref{sec:setup}. Assume the actual-fiber criterion \eqref{eq:fiber-ratio}, or an independently verified comparison \eqref{eq:task-compare}; the actual conditional mass bounds \eqref{eq:mass-upper}--\eqref{eq:mass-lower}; the lower effort certificate \eqref{eq:effort-drift}; and the $B$-live-label program $\mathcal G$. Set
\[
 b_-(N)=uN/\kappa_-,\qquad
 b_+(N)=\max\{1,\ell N/\kappa_+\}.
\]
For $N\ge\kappa_+$ and $m\ge1$,
\begin{equation}\label{eq:cp-main}
 k_-\Qm_m(b_-(N))\le R_\cp^\Theta(N,m)-B_0
       \le k_+\frac u\ell\Qm_m(b_+(N)).
\end{equation}
For $M\ge B+1$,
\begin{equation}\label{eq:aut-main}
 k_-\Qm_{M-1}(b_-(N))\le R_\aut^\Theta(N,M)-B_0
       \le k_+\frac u\ell\Qm_{M-B}(b_+(N)).
\end{equation}
Every finite upper is implemented at the original budget by one common Borel routing program with exactly the stated upper alphabet bound. Proposition~\ref{prop:finite-cells} supplies a finite-description implementation under its explicit computational certificate. No exact finite observable primary driver is required. No interchange of worst-model and best-program operations is used.

If every legal autonomous procedure requires at least $b_0$ live labels under one common admitted model, the lower profile in \eqref{eq:aut-main} is $\Qm_{M-b_0}$. In particular, if $b_0=B$, the same paid terminal width occurs on both sides. Section~\ref{sec:memory} supplies a physical-command certificate for this condition. No procedure is feasible for $N<\kappa_-$. All profile minima at finite $b\ge1$ exist.
\end{theorem}

The two resource arguments are intentional. Bounded conditional density ratios do not imply constant-factor comparison of $\Qm(cb)$ and $\Qm(b)$ at a critical atom. Theorem~\ref{thm:atomic} preserves that transition. Under regular acquired geometry, Corollary~\ref{cor:regular} gives a matching law. Thus the theorem supplies both an all-strategy lower and a finite-state upper, rather than identifying a relaxation with an implementation by definition.

A physical experiment may contain several additive resources. The same construction bounds the expected amount of resource $j$ by $\kappa_{+,j}b/\ell$ whenever the acquisition program has that conditional attempt-cost bound. To obtain an upper satisfying all budgets, choose $b$ using their minimum admissible value. A lower can use any nonnegative linear combination of resources with a verified drift inequality. One constraint with a positive per-call lower bound supplies finite expected duration; bytes, raw calls and physical time are not summed without explicit conversion weights.

\begin{corollary}[Equivalence of certified task marks]\label{cor:mark-invariance}
If two bounded marks $(f,w)$ and $(\widetilde f,\widetilde w)$ on the same observed mark space satisfy uniform comparisons with the same full conditional task, their marked distortions are comparable for every prediction, and their acquired profiles are comparable with the same $b,m$. In particular any such mark can replace \eqref{eq:canonical-mark}, with explicit changes only in the comparison constants.
\end{corollary}
\begin{proof}
For each observed $x$ choose any attainable conditional pair on its nonempty task fiber and compare both marked distortions to its risk. The inequalities hold for every $r$, hence also for the minimum over a common finite dictionary. Integration against any common $h\nu$ and then minimization prove the profile comparison. No change of preparation or renormalization of a stratum is involved.
\end{proof}

\section{Proof of the transfer theorem}
\begin{lemma}[Predictable effort and terminal domination]\label{lem:domination}
Every finite-cost legal procedure satisfies
\[
 \kappa_-\E_\theta\sum_{t\le\tau}A_t
 \le\E_\theta\sum_{t\le\tau}c_t.
\]
Its terminal mark law $\lambda_\theta$ obeys
\begin{equation}\label{eq:terminal-domination}
 \lambda_\theta\le u\left(\E_\theta\sum_{t\le\tau}A_t\right)\nu.
\end{equation}
These conclusions allow arbitrary history-dependent commands and stopping using the full report.
\end{lemma}
\begin{proof}
Because $\{t\le\tau\}$ is measurable before call $t$, sum \eqref{eq:effort-drift} to $\tau\wedge n$. The expectation of the terminal potential tends to zero: it is zero on $\{\tau\le n\}$, uniformly bounded elsewhere, and $\tau<\infty$ almost surely. Monotone convergence for the nonnegative cost and count yields the first assertion. At least one productive call is necessary, so $N<\kappa_-$ is infeasible.

For a Borel set $D\subset\mathsf X$, the actual terminal subcount is bounded pathwise by the full productive count:
\[
 \one_{\{X_\tau\in D\}}
 \le\sum_{t\le\tau}A_t\one_{\{X_t\in D\}}.
\]
Take expectations, condition before each productive report, and use \eqref{eq:mass-upper}. This proves \eqref{eq:terminal-domination} first as a measure inequality. Radon--Nikodym then gives $\lambda_\theta=h_\theta\nu$ with $0\le h_\theta\le uN/\kappa_-$ and mass one. Informative selection through $V_t$ cannot invalidate domination by the full acquired count. This step does not assert that the acquired count is a constant multiple of $\nu$.
\end{proof}

\begin{lemma}[Correlated rejection with a paid acquisition program]\label{lem:rejection}
Fix $b\ge1$, a dictionary $C$ and a Borel density $0\le h\le b$ of mass one. Run $\mathcal G$ repeatedly. At a productive mark $x$, accept with probability $h(x)/b$, using a fresh coin; on acceptance output a nearest dictionary value, and on rejection return to its designated live label. Then the number $J$ of attempts has $\E_\theta J\le b/\ell$, the expected raw cost is at most $\kappa_+b/\ell$, and the terminal marked distortion is at most
\[
 \frac u\ell\int d_C h\,d\nu.
\]
The construction uses $B+m$ labels. It does not assume independent attempts.
\end{lemma}
\begin{proof}
At every productive history the acceptance probability is at least $\ell/b$, by \eqref{eq:mass-lower}. The same is true conditional on the beginning of its attempt, by the tower property. Thus
\[
 \Pp_\theta(J>j)\le(1-\ell/b)^j,
 \qquad \E_\theta J\le b/\ell.
\]
Let $C_j$ be the raw cost of the $j$th attempt when started. The event $\{J\ge j\}$ is known at its start, and its conditional mean cost is at most $\kappa_+$. Tonelli and conditioning give
\[
 \E_\theta\sum_{j\le J}C_j
 =\sum_{j\ge1}\E_\theta[\one_{\{J\ge j\}}\E_\theta(C_j\mid\text{attempt start})]
 \le\kappa_+\E_\theta J.
\]
No independence between $C_j$, the report and acceptance was used.

Conditional on each productive history, the expected distortion contributed by acceptance is at most $u b^{-1}\int d_C h\,d\nu$. Summing over actual attempts gives the claimed terminal bound. There is exactly one terminal contribution almost surely. The extra report is ignored only by this particular upper construction, not by the competing lower class. The live labels execute $\mathcal G$ and return on rejection; the $m$ terminal labels store the fixed output. Private randomness is fresh and leaves no uncharged seed.
\end{proof}

\begin{proof}[Proof of Theorem~\ref{thm:main}]
For every procedure and every model, apply \eqref{eq:task-compare} to its actual prediction. Randomizing over a fixed squared-loss dictionary cannot improve upon its nearest member. Lemma~\ref{lem:domination} gives the lower profile for its actual terminal measure. Take the worst model and then the infimum over common procedures. An autonomous procedure has at least one live label, hence at most $M-1$ fixed terminal values. If $b_0$ live labels are mandatory under a common admitted model, use that model for the sharpened lower; the other modelwise inequalities still hold. The common-model quantifier is important.

For attainment in \eqref{eq:profile}, the densities form a weak-star compact subset of $L^\infty(\nu)$: $0\le h\le b$ is closed by nonnegative $L^1$ tests and the mass constraint is closed. Use compactness with nets; no sequential compactness of general stopped trajectory laws is asserted. The dictionary lies in compact $H^m$. Its squared-distance function converges uniformly when its centers converge, since $H$ and $f$ are bounded. The objective is therefore jointly continuous and a minimum exists. Choose Borel representatives clipped into $[0,b]$ and a deterministic Borel tie rule for nearest centers. Such changes on $\nu$-null sets do not change any actual law because of \eqref{eq:mass-upper}.

If $\ell N/\kappa_+>1$, apply Lemma~\ref{lem:rejection} at $b=\ell N/\kappa_+$ to a minimizer. Its cost is at most $N$ in every model; multiply its distortion bound by $k_+$. If $\ell N/\kappa_+\le1$, take $b=1$. The only feasible density is $h=1$ almost everywhere, so acceptance is certain at the first productive report. Its cost is at most $\kappa_+\le N$, and its distortion at most $u\Qm_m(1)$, which is bounded by the displayed $u/\ell$ factor. This endpoint argument avoids falsely charging $\kappa_+/\ell$ for certain acceptance. The upper also qualifies as a checkpoint procedure. Setting $m=M-B$ gives the autonomous upper.
\end{proof}

\begin{example}[Why exact primary closure cannot be inferred]\label{ex:closure}
Let a hidden sign have two different posterior probabilities after histories $h_1,h_2$ in the same visible gate phase, and let its next success probabilities be $p_+$ and $p_-$ conditional on the sign. Then
\[
 \Pp(\mathrm{success}\mid h_i)
   =p_-+(p_+-p_-)\Pp(E=+1\mid h_i)
\]
differs with $i$. Multiplying a total phase count by one fixed success row is wrong. Nevertheless each conditional probability lies in $[p_-,p_+]$ and the effort and program certificates of Theorem~\ref{thm:herald} hold. The theorem does not repair the false equality by replacing the true history with an averaged artificial state.
\end{example}

\section{When a history-uniform task mark exists}
\label{sec:fiber}
\begin{theorem}[Conditional minimax radius criterion]\label{thm:fiber}
Under the fiber regularity in Section~\ref{sec:setup}, \eqref{eq:canonical-mark} defines bounded Borel functions with a unique center. The following are equivalent, up to the displayed constants.

\smallskip\noindent
(i) There are Borel $f,w\ge0$ and $0<a\le A<\infty$ such that, for every $(\mu,s)\in\mathcal F_x$ and $r\in H$,
\[
 a\{w+\|f-r\|^2\}\le s+\|\mu-r\|^2
       \le A\{w+\|f-r\|^2\}.
\]
\noindent
(ii) $w_*(x)\le K\underline s(x)$ uniformly for some finite $K$.

Condition (i) implies (ii) with $K=A/a$. Condition (ii) implies (i) for the canonical mark with $a=(3K)^{-1}$, $A=2$. Zero-radius fibers are included without division by zero. This is a criterion for the stated all-prediction, all-history quadratic comparison, not a necessary condition for every finite-resource risk theorem.
\end{theorem}
\begin{proof}
Each function $r\mapsto s+\|\mu-r\|^2$ is strongly convex with the same quadratic part. Hence
\[
 g_x(\lambda r+(1-\lambda)r')
 \le\lambda g_x(r)+(1-\lambda)g_x(r')
       -\lambda(1-\lambda)\|r-r'\|^2.
\]
It is uniformly Lipschitz on compact $H$, uniformly over $x$. Existence follows by compactness and uniqueness by this inequality. To see Borel measurability without an implicit selection assumption, minimize $g_x$ over successive fixed finite nets of a countable dense subset of $H$, breaking ties by index. These minimizers are Borel. Their objective errors tend uniformly to zero by Lipschitz continuity, and strong convexity on the segment to the unique minimizer implies convergence to it. Therefore $f_*$ and $w_*$ are Borel.

Under (i), putting $r=f$ gives $g_x(f)\le A w$. Putting $r=\mu$, separately for every pair in the fiber, gives $aw\le s$. Thus
\[
 w_*\le g_x(f)\le Aw\le(A/a)\underline s.
\]
The quantifier over every fixed $r\in H$ permits this pointwise substitution; it does not give the observer access to $\mu$.

Conversely let (ii) hold. For every conditional pair,
\[
 s+\|\mu-f_*\|^2\le w_*,\qquad s\ge w_*/K.
\]
The upper comparison follows from
\[
 s+\|\mu-r\|^2
 \le2\{s+\|\mu-f_*\|^2\}+2\|f_*-r\|^2
 \le2\{w_*+\|f_*-r\|^2\}.
\]
For the lower,
\[
 w_*+\|f_*-r\|^2
 \le3w_*+2\|\mu-r\|^2
 \le3K\{s+\|\mu-r\|^2\}.
\]
If $w_*=0$, all pairs have $s=0$ and $\mu=f_*$, and the comparison is immediate. Also $w_*\ge\underline s$, so $K\ge1$ can always be used on positive fibers.
\end{proof}

\subsection*{What the criterion distinguishes}
An exact task identity has a singleton fiber $(f,w)$ and directly yields the mark. A primitive minorization or positivity-improving channel can instead bound all conditional variances below and conditional means above, producing a sufficient certificate. A calibration supplied by an external protocol must pay for that protocol when it forms part of the experiment. None of these cases follows from positivity of an arbitrary normalized kernel alone.

\begin{example}[Repeated information can destroy a current-report mark]\label{ex:calibration-failure}
Let $E\in\{-1,1\}$ be fixed, initially fair, and let independent reports satisfy $\Pp(V=E)=1-\gamma>1/2$. Set the audit to $Z=E$. Repeated positive reports drive $\Pp(E=-1\mid V_1=\cdots=V_n=1)$ to zero. Histories ending in the same current report $V_n=1$ also include histories with a substantial negative posterior. Their fiber has least variance zero but positive minimax radius. Thus no positive uniform quadratic current-report mark exists. Enlarging the predictive state to retain the posterior is a different, separately charged implementation; the failed mark cannot be rescued by declaring its formal support acquired.
\end{example}

\begin{proposition}[Necessary dependence on a vanishing positivity floor]\label{prop:q-sharp}
For the binary audit $Z=y+tE$, suppose a pre-report preparation can have sign probabilities $(1-q,q)$ and a weak sign report has likelihoods $1-\gamma,\gamma$, with $0<\gamma<1/2$. On a positive report the smaller posterior probability is
\[
 a_q=\frac{q\gamma}{(1-q)(1-\gamma)+q\gamma}.
\]
Its residual variance is $4a_q(1-a_q)t^2=O(qt^2)$. If the admitted family also contains the sign-reversed preparations, its conditional minimax radius is of order $t^2$ for small $q$. Consequently the best uniform ratio in Theorem~\ref{thm:fiber} is at least of order $q^{-1}$. The loss of uniformity as $q\downarrow0$ in this mark mechanism is necessary.
\end{proposition}
\begin{proof}
Bayes' rule gives the displayed fraction. The two extreme preparations, conditional on the same weak positive report, produce audit means tending to $y+t$ and $y-t$. Any center has squared distance at least a fixed multiple of $t^2$ from one of them. The least variance is bounded by the first preparation's variance. Their ratio has the stated order. The same calculation is obtained from diagonal qubit states and a weak $\sigma_z$ effect. It is a subfamily of an instrument class that may also contain noncommuting commands. This proves a limitation of uniform proxy calibration, not a complete minimax law in the joint $q,N,M$ limit.
\end{proof}

\section{Physical commands and live-state width}
\label{sec:memory}
The observer's action at a live label is chosen before the next report. A command means its physical instrument, its port of application, its legal domain, its recorded outputs and its cost, not a string assigned to it. Two command names are aliases when they represent this same object on every relevant raw history. A \emph{stateless physical recoding} replaces names and report coordinates by fixed measurable translations preserving these objects, without reading a hidden phase or using additional memory. A translation that must recover a protocol phase is a simulator, whose state is separately charged.

Let $I$ be a finite collection of mandatory contexts. Context membership is a property of the raw physical history, not information freely supplied to a label-only observer. For each $i\in I$, let $\mathcal U_i$ be its set of legal physical commands; it is the same set at every history of that context. Define
\begin{equation}\label{eq:semantic-width}
 b_{\rm phys}(I)=\min\left\{r:I=I_1\cup\cdots\cup I_r,
       \ \bigcap_{i\in I_j}\mathcal U_i\ne\varnothing\ (1\le j\le r)\right\}.
\end{equation}
Empty pieces are omitted. This is a compatibility cover, not merely a count of different action names. The command space and its legal sets are Borel. For the lower below, contexts sharing a label must put probability one on the intersection of their legal command sets.

\begin{proposition}[Semantic mandatory-live-state lower]\label{prop:semantic}
Suppose that under one admitted model $\theta_0$, every legal terminating procedure visits every context in $I$ with positive probability. Every autonomous observer then has at least $b_{\rm phys}(I)$ live labels. This lower is unchanged under stateless physical recodings and duplication of aliases. In particular, pairwise disjoint physical command sets of $B$ mandatory contexts imply $B$ paid live labels.
\end{proposition}
\begin{proof}
For a particular observer and $\theta_0$, assign to each live label $s$ the set $I_s$ of contexts in which it occurs with positive probability. Its pre-report command distribution is fixed by $s$. Exact legality implies that this distribution gives probability one to every $\mathcal U_i$, $i\in I_s$. Since $I$ is finite, it gives probability one to their intersection, which is therefore nonempty. Positive-probability visits imply that the sets $I_s$ cover $I$. Their number is at least \eqref{eq:semantic-width}. The argument uses the same model for all contexts and does not compare separate procedures chosen for separate models. A stateless recoding transports exactly the legal physical command intersections, so it preserves the cover number. A duplicate name adds no physical command and changes no intersection.
\end{proof}

\begin{example}[Aliasing and a genuinely stronger interface]\label{ex:alias}
Suppose a physical ``advance'' operation is legal in two phases, with the same command available in both. Writing it as two phase-decorated aliases does not make the physical legal sets disjoint. Counting these names would give a false two-live-label lower. Conversely, if a gate pulse and an acquisition pulse are physically distinct and each is forbidden at the other phase, their sets are disjoint. Adding a new ``advance'' instruction that detects the phase and selects the correct pulse enlarges the interface. If implemented externally, its phase detector is a simulator, not a free renaming. Proposition~\ref{prop:semantic} is intrinsic under the stated resource-preserving presentations; it is not an invariant under adding physical feedback or uncharged external state.
\end{example}

The cover number is a lower certificate, not a characterization of the entire minimal protocol automaton. Routing compatibility can require more states than first-command compatibility. The matching applications below have two mandatory physical ports and an explicit two-label routing program, so this distinction leaves no state gap there. Persistent labels still do not equal total bit complexity: a program table, arithmetic workspace, a retained random seed and a clock are additional resources whenever present.

\section{Acquired geometry, singular mass and state-dependent waiting}
\label{sec:geometry}
The following quantization calculation is inherited from the stopped-experiment analysis \cite{GTF52,GTF53}; we include its proof, constants' range and singular extension to keep the present transfer self-contained. Put
\[
 \sigma=(f,\sqrt w)_\#\nu.
\]
Assume its location support is a compact $K\subset H$ and $0\le\sqrt w\le1$. Fix $d,\alpha>0$, $r_0\in(0,1]$, and $c,C>0$ such that for every $x\in K$, $0<r\le r_0$ and $0<s\le1$,
\begin{equation}\label{eq:rectangle}
 cr^d s^\alpha\le\sigma\bigl(B(x,r)\times[0,s]\bigr)
                  \le Cr^d s^\alpha.
\end{equation}
Thus the geometry is that of the actual marked law. In particular it can be singular in its ambient Euclidean space. A fixed finite cover of $K$ by radius $r_0$ balls gives $\sigma(H\times[0,s])\le C_Ts^\alpha$.

\begin{proposition}[Marked selection and quantization]\label{prop:geometry}
Under \eqref{eq:rectangle}, uniformly for $b\ge1$, $m\ge1$,
\begin{equation}\label{eq:regular-profile}
 \Qm_m(b)\asymp
 F(b,m):=\max\{b^{-2/\alpha},(bm)^{-2/(d+\alpha)}\}.
\end{equation}
Constants depend only on the displayed geometric data and the diameter of $H$, not on $b,m$.
\end{proposition}
\begin{proof}
A mass-one density at most $b$ cannot put more than $bC_Ts^\alpha$ mass at reliability at most $s$. Taking $s=(2C_Tb)^{-1/\alpha}$, reduced by a fixed constant if necessary, gives a lower of order $b^{-2/\alpha}$.

For any dictionary of $m$ centers in $H$, a marked distortion sublevel $d_C\le t^2$ lies in the union of $m$ rectangles with reliability at most $t$ and location radius $t$. A center not in $K$ can be moved for this estimate to a point of $K$ within $t$, expanding the radius to $2t$. For $2t\le r_0$ its mass is at most $C_Rm t^{d+\alpha}$. Set $t=c_0(bm)^{-1/(d+\alpha)}$, with
\[
 c_0\le \min\{r_0/2,1,(2C_R)^{-1/(d+\alpha)}\}.
\]
At least half the selected mass has distortion greater than $t^2$. This proves the other lower for the entire range $b,m\ge1$.

The marginal at $s=1$ has two-sided ball mass of order $r^d$. A maximal separated-set argument supplies, for every integer $k\ge1$, $k$ points in $K$ separated by at least $a k^{-1/d}$, with one fixed $a>0$. For completeness, a maximal $r$-separated set covers $K$ by $r$ balls, so its cardinality is at least $(Cr^d)^{-1}$ for $r\le r_0$; choose $r$ a sufficiently small multiple of $k^{-1/d}$ and retain $k$ points. The $k=1$ case imposes no separation condition.

Choose $C_1$ so that $cC_1^{d+\alpha}\ge1$, then choose $c_1\in(0,1]$ so small that
\[
 C_1 c_1^{\alpha/[d(d+\alpha)]}\le a/3.
\]
Increase the explicit threshold
\begin{equation}\label{eq:b-threshold}
 b_0=\max\left\{1,(2/c_1)^{\alpha/d},
        (C_1/\min\{r_0,1\})^{d+\alpha}\right\}
\end{equation}
if needed to accommodate the chosen packing constant. For $b\ge b_0$ set
\[
 k=\min\{m,\lfloor c_1 b^{d/\alpha}\rfloor\},\qquad
 t=C_1(bk)^{-1/(d+\alpha)}.
\]
The radius satisfies $t\le\min(r_0,1)$ and $3t\le a k^{-1/d}$. The $k$ reliability rectangles are therefore disjoint and their union has mass at least $kc t^{d+\alpha}\ge1/b$. On this union the dictionary formed by the $k$ centers has distortion at most $2t^2$. Normalize its restriction to mass one; its density is at most $b$. If $k=m$ this is the second term of $F$; otherwise $k$ is bounded below by a fixed multiple of $b^{d/\alpha}$ and it is the first term.

For $1\le b<b_0$, use density one and a single center. Its distortion is at most $1+\operatorname{diam}(H)^2$, while $F(b,m)\ge b_0^{-2/\alpha}$. This handles every $m$, including arbitrarily large $m$; it is a bounded-budget argument, not an assertion that the remaining $(b,m)$ set is compact.
\end{proof}

\begin{corollary}[Matching curves under regular acquired geometry]\label{cor:regular}
Assume Theorem~\ref{thm:main}, \eqref{eq:rectangle}, and $\kappa_+/\kappa_-\le\Gamma<\infty$. Then, for $N\ge\kappa_+$,
\[
 R_\cp^\Theta(N,m)-B_0\asymp F(N/\kappa_+,m).
\]
If $B$ mandatory physical contexts verify Proposition~\ref{prop:semantic} and the upper uses $B$ live labels, then for every $M\ge B+1$,
\begin{equation}\label{eq:matched-aut}
 R_\aut^\Theta(N,M)-B_0\asymp F(N/\kappa_+,M-B).
\end{equation}
The constants additionally depend on $\ell,u,\Gamma,k_-,k_+$, but are independent of $N,M$ and of the absolute scale of $\kappa_+$.
\end{corollary}
\begin{proof}
Apply Proposition~\ref{prop:geometry} to the two arguments of Theorem~\ref{thm:main}. Multiplying $b\ge1$ by fixed factors changes either power by fixed factors. The upper endpoint $\max(1,\ell b)$ obeys the same comparison. Mandatory width gives the same $M-B$ on both sides. No regular variation in an ambient parameter dimension is used.
\end{proof}

\begin{theorem}[An actual atomic transition]\label{thm:atomic}
Let the marked acquired law be
\[
 \nu=\lambda\delta_{(x_*,0)}+(1-\lambda)\nu_0,
 \qquad 0<\lambda<1,
\]
where $\nu_0$ is a probability law satisfying \eqref{eq:rectangle}; the notation here identifies marks with their location and reliability. Write $r=(1-b\lambda)_+$. If $r=0$, then $\Qm_m^\nu(b)=0$ for $m\ge1$. If $r>0$ and $m\ge2$, then
\begin{equation}\label{eq:atomic-profile}
 r\Qm_m^{\nu_0}\!\left(\frac{b(1-\lambda)}r\right)
 \le\Qm_m^\nu(b)
 \le r\Qm_{m-1}^{\nu_0}\!\left(\frac{b(1-\lambda)}r\right).
\end{equation}
Together with \eqref{eq:regular-profile}, this is a matching curve, including its vanishing at the actual-mass threshold $b\lambda=1$.
\end{theorem}
\begin{proof}
The atomic selected mass is at most $b\lambda$, so the regular part has mass at least $r$. Removing surplus regular mass can only decrease its nonnegative distortion. A regular measure of mass exactly $r$, normalized to probability, has density at most $b(1-\lambda)/r$ relative to $\nu_0$. This proves the lower with the original dictionary. For the upper put selected mass $b\lambda$ on the atom, use one center at $x_*$, and use the other $m-1$ centers for an optimal regular submeasure of mass $r$. If $b\lambda\ge1$, select the atom alone with density $1/\lambda\le b$. For $m\ge2$, $m-1$ and $m$ differ by at most a factor two, proving matching through Proposition~\ref{prop:geometry}. The factor $r$ and the threshold $b\lambda$ cannot be removed by renormalizing the regular stratum.
\end{proof}

For example, a product of an Ahlfors-regular Cantor location law and the reliability law $\Pp(T\le s)=s^\alpha$ satisfies \eqref{eq:rectangle}. Products with unequal factor dimensions give anisotropic cell constructions before the scalar distortion is chosen. An actual revealing atom can be supplied as an additional reported outcome. These are singular acquired laws, not smooth parameter-space surrogates. Under conditional density envelopes the two sides of Theorem~\ref{thm:main} retain their respective atomic thresholds; without additional exact-law structure, they must not be collapsed into one threshold.

\begin{theorem}[Unknown, history-dependent heralding]\label{thm:herald}
There are two physical phases, gate and acquisition. A gate call costs one and permits no stopping. Its conditional success probability $p_h$, allowed to depend on the entire history, command and unknown fixed model, satisfies
\[
 0<p_-\le p_h\le p_+\le1.
\]
Failure leaves the gate phase; success enables one acquisition call. The latter costs one, is productive, and returns to the gate phase on rejection. Only the corresponding physical command is legal at each of the two ports; both ports are mandatory under one admitted model. Assume the task comparison and conditional mark envelopes of Section~\ref{sec:setup}. Then Theorem~\ref{thm:main} applies with
\begin{equation}\label{eq:kappa}
 B=2,\qquad \kappa_-=1+1/p_+,\qquad \kappa_+=1+1/p_-.
\end{equation}
In particular \eqref{eq:matched-aut} holds under \eqref{eq:rectangle} and $p_+/p_-\le\Gamma$. The program needs only the certified bounds, not the true $p_h$. The hidden system is not reset between attempts.

No procedure is feasible for $N<\kappa_-$. If the family contains a model with gate success constantly $p_-$ and the same mandatory protocol, robust feasibility is exactly $N\ge\kappa_+$. Uniformity as $p_-\downarrow0$ is on this rescaled feasible budget. At fixed finite $N$, that family's problem becomes infeasible when $1+1/p_->N$. At $p_-=p_+=1$ both one gate and one acquisition call are paid, so $\kappa_+=2$.
\end{theorem}
\begin{proof}
For the all-strategy effort lower, put $V=0$ in gate phase and $V=-1/p_+$ in acquisition phase. After every productive call, committed or not, give it value zero. On a gate call,
\[
 \E[1+V_{\rm next}-V_{\rm current}\mid h]=1-p_h/p_+\ge0.
\]
On an acquisition call the increment is $1+1/p_+=\kappa_-$. These are precisely \eqref{eq:effort-drift}, without a fixed transition row.

For an attempt program, repeat the gate command until success and then acquire. The two-label cost supersolution is $W_G=\kappa_+$, $W_A=1$. Indeed
\[
 1+(1-p_h)W_G+p_hW_A\le W_G,\qquad 1\le W_A.
\]
Lemma~\ref{lem:navigator} proves the uniform conditional attempt bound. Proposition~\ref{prop:semantic} gives two live labels in the lower, since these are different physical ports, not aliases. Thus all hypotheses of Theorem~\ref{thm:main} have been checked from conditional raw laws. When the constant-$p_-$ model belongs to the family, predictable gate counts give expected cost at least $\kappa_+$ for one productive call; accepting that first call attains it. This proves the exact feasibility statement. No independent-cycle or true-rate argument appears.
\end{proof}

\section{Two raw-kernel verifications with endogenous primary laws}
\label{sec:realizations}
Both classes below use the same physical gate/acquisition protocol, but their primary report laws depend on their hidden conditional states. Consequently neither has the exact common finite primary driver of \cite{GTF53}. Let $\nu_0(dy,dt)$ be a specified probability law on compact $K\times[0,1]$, let $|\phi_a|\le1$ with $\int\phi_a\,d\nu_0=0$, and fix $0\le\beta<1$. Put
\[
 c=1-\beta,\quad C=1+\beta,\quad
 0<q\le\tfrac12,\quad 0<\gamma<\tfrac12,\quad
 a_0=\frac{qc\gamma}{C(1-\gamma)}\le\tfrac12.
\]
The functions, bounds and a reference model are known; the actual transition, gate and weak-observation parameters need not be known. All admitted actions obey the same bounds. The audit carrier is chosen compact and convex to contain $y+t e\,e_1$ for $y\in K$, $0\le t\le1$ and $e=\pm1$.

\begin{theorem}[Controlled hidden-state sensing]\label{thm:classical}
Let a hidden sign evolve on every call by a controlled stochastic matrix whose two entries in each row lie in $[q,1-q]$. On a gate call its state-dependent success probabilities lie in $[p_-,p_+]$. On an acquisition call, after that transition, the primary mark has conditional density
\[
 L_a(y,t\mid e)=1+\beta e\phi_a(y,t)
\]
relative to $\nu_0$. An auxiliary weak sign report $v\in\{-1,1\}$ is then produced with each likelihood in $[\gamma,1-\gamma]$. The audit is
\[
 Z=Y+TEe_1,
\]
with the current hidden sign and no extra transition. Then $f(y,t)=y$, $w(y,t)=t^2$, $B_0=0$ satisfy \eqref{eq:ordinary-calibration} with
\begin{equation}\label{eq:raw-calibration}
 a_v=4a_0(1-a_0),\qquad A_v=a_b=1.
\end{equation}
An actually acquired reference law $\nu$ gives conditional envelope constants $\ell=c/C$, $u=C/c$. Thus Theorems~\ref{thm:main} and \ref{thm:herald} apply to every legal history-dependent competitor and one common two-label upper program. If $\nu_0$ has the marked rectangle geometry, Corollary~\ref{cor:regular} gives its matching curve with the same $d,\alpha$.
\end{theorem}
\begin{proof}
Immediately before the primary acquisition report, both sign probabilities are at least $q$, regardless of how much information was obtained or rejected earlier. Since $c\le L_a\le C$, the primary mark's conditional density lies in $[c,C]$ relative to $\nu_0$. After also observing $v$, the numerator of \emph{each} sign posterior is at least $qc\gamma$, while the report likelihood denominator is at most $C(1-\gamma)$. Both posteriors therefore lie in $[a_0,1-a_0]$. In particular $a_0\le1/2$, and if $b=\E[E\mid h,y,t,v]$, then
\[
 \operatorname{Var}(Z\mid h,y,t,v)=(1-b^2)t^2
        \in[4a_0(1-a_0)t^2,t^2],\qquad
 \|\E[Z\mid h,y,t,v]-y\|^2=b^2t^2\le t^2.
\]
This is \eqref{eq:raw-calibration} after the entire report. A fresh stopping coin independent of the audit does not change the conditional identity.

Run the declared gate/acquisition protocol once in a known reference member of the family. Its first productive mark law $\nu$ is an actual acquisition law; conditional densities and their mixtures show $c\nu_0\le\nu\le C\nu_0$. Hence every conditional productive law in every model is between $(c/C)\nu$ and $(C/c)\nu$. This comparison does not supply the true model to the upper program or condition away a tail event. It also transfers rectangle estimates from $\nu_0$ to $\nu$.

The raw hidden transition may have different rows, and gate and primary likelihoods may depend on the sign. Its posterior changes after failed gates and auxiliary reports. In that case both future gate probabilities and primary mark densities change with the observed history, as in Example~\ref{ex:closure}. All preceding estimates remain valid. The program stores its physical phase, not this varying posterior. Normalization denominators on acquisition are bounded below by $c\gamma$ relative to $\nu_0$; thus the report and posterior maps have explicit uniform continuity bounds on the compact pre-report simplex. No contraction or resetting of the hidden process was used in the transfer proof.
\end{proof}

\begin{theorem}[Noncommuting quantum heralding]\label{thm:quantum}
Let the hidden system be a qubit. On each call first apply
\[
 \mathcal D_q(\rho)=(1-2q)\rho+q\tr(\rho)I,
\]
then a controlled unitary $U_a$. Gate effects $G_a,I-G_a$ satisfy $p_-I\le G_a\le p_+I$ and use their positive square roots as Kraus operators. On an acquisition call the primary instrument density is $L_a(y,t)=F_a(y,t)^{1/2}$, where
\[
 F_a(y,t)=I+\beta\phi_a(y,t)n_a\cdot\sigma,
 \qquad \int F_a\,d\nu_0=I.
\]
It is followed by weak effects $E_{a,v}=(I+v\eta m_a\cdot\sigma)/2$, $\eta=1-2\gamma$, and Kraus operators $K_{a,v}=E_{a,v}^{1/2}$. The axes need not commute. The audit is $Z=Y+Te\,e_1$, with $e$ the outcome of a subsequent $\sigma_z$ measurement and no intervening channel. Then the same constants \eqref{eq:raw-calibration} and the same conditional envelopes verify Theorems~\ref{thm:main} and \ref{thm:herald}.
\end{theorem}
\begin{proof}
For a normalized input, the state after channel and unitary satisfies $\rho'\ge qI$. The primary effect obeys $cI\le F_a\le CI$. Its density $\tr(F_a\rho')$ lies in $[c,C]$, and the integrated instrument is trace preserving. After the primary and weak reports the unnormalized state is
\[
 \omega=K_{a,v}L_a\rho' L_aK_{a,v}\ge qc\gamma I.
\]
For normalization,
\[
 \tr\omega=\tr(E_{a,v}L_a\rho' L_a)
       \le(1-\gamma)\tr(F_a\rho')\le C(1-\gamma).
\]
Thus $\omega/\tr\omega\ge a_0I$. Each $\sigma_z$ outcome has probability at least $a_0$, so its mean and variance give precisely the classical audit comparison. The reference acquisition and density envelopes are proved in exactly the same measure-theoretic way, now using $\tr(F_a\rho')\in[c,C]$. Gate success $\tr(G_a\rho')\in[p_-,p_+]$ is generally history dependent. Its state update is not a reset.

Nonparallel $n_a,m_a$ and gate axes yield noncommuting instruments; no shared basis was used. The physical state enters the raw experiment but is not writable or freely readable controller memory. The observer remains a two-live-label classical program. Report-kernel continuity follows from linearity of unnormalized instruments and the positive normalization bounds. All arguments are finite-dimensional. They assert neither unbounded-operator closure nor uniform control at projective or vanishing-$q$ limits.
\end{proof}

\begin{lemma}[Raw report continuity and normalized update stability]\label{lem:raw-continuity}
In either realization, at a productive call the complete primary/weak report law is nonexpansive in total variation with respect to the input state distance, up to the conventional factor $1/2$. For a fixed positive report its normalized predictive update satisfies
\[
 \|F_{a,x,v}(z)-F_{a,x,v}(z')\|_1
 \le \frac{2C(1-\gamma)}{c\gamma}\|z-z'\|_1,
\]
where the norm is total coordinate variation classically and trace norm quantum mechanically. This bound is uniform over the declared acquisition family. It is a one-step stability estimate, not a contraction assertion for arbitrary long products.
\end{lemma}
\begin{proof}
Write the unnormalized report map as $A$. Positivity gives
$\|A(z-z')\|_1\le C(1-\gamma)\|z-z'\|_1$ by the positive/negative decomposition (or its finite-dimensional matrix analogue). Its trace or total mass is at least $c\gamma$ on normalized inputs. Subtract $Az/\tr Az$ and $Az'/\tr Az'$; one term is the unnormalized difference divided by the first denominator, and the other is bounded by the same quantity because the second normalized output has norm one. This proves the displayed factor two. Integrating the absolute report-mass differences uses the trace-preserving complete instrument and gives the total-variation assertion. The preceding transition/channel and unitary are nonexpansive in these norms. No derivative of a posterior at a zero-evidence event is taken.
\end{proof}

\subsection*{Singular marks and the sharp positivity boundary}
Take $\nu_0$ on a Cantor location support with reliability law $\Pp(T\le s)=s^\alpha$. Bounded likelihood factors preserve its actual singular support and rectangle masses; the two raw proofs require no Lebesgue density in ambient coordinates. Alternatively add a marked atom $(y_*,0)$ to the preparation/report instrument, preserving its actual mixture weight in the reference law. Theorems~\ref{thm:main} and \ref{thm:atomic} then retain the revealing transition. The statement is not uniform under arbitrarily changing supports.

The explicit dependence on positivity is $a_0=qc\gamma/[C(1-\gamma)]$, followed by $a_v=4a_0(1-a_0)$ and the constants in \eqref{eq:ordinary-calibration}. Proposition~\ref{prop:q-sharp} shows that this deterioration is necessary for a current-mark calibration over the full admitted family. Uniformity in rare gate success with bounded $p_+/p_-$ must therefore not be confused with uniformity in $q\downarrow0$ or $\gamma\downarrow0$.

\subsection*{A conventional Gaussian audit without a supplied reliability coordinate}
The preceding marked models make the geometry transparent. The next consequence uses the ordinary task of predicting the current hidden sign; its reliability is derived from the raw likelihood, not supplied as an additional random coordinate. Its Gaussian tail mechanism is inherited from the stopped-experiment companion~\cite{GTF52}; the new application pays for history-dependent heralding through the same general theorem.

\begin{theorem}[Endogenously gated Gaussian prediction]\label{thm:gaussian-herald}
Use the two physical ports of Theorem~\ref{thm:herald}. A hidden sign evolves on every call with controlled transition entries in $[q,1-q]$, $0<q\le1/2$. Gate probabilities may depend on its hidden state and lie in $[p_-,p_+]$, with $p_+/p_-\le\Gamma$. A productive call reports
\[
 Y=aE+\sigma G,\qquad G\sim N(0,1),\qquad a,\sigma>0,
\]
with fresh $G$, and the subsequent audit is simply $Z=E$, without an extra transition. The transition and gate parameters are inaccessible. Suppose the known reference family includes a fair post-transition sign and a sign-independent gate. Put $B_0=0$, $\kappa_+=1+1/p_-$ and
\[
 F_G(b)=\exp\!\left[-\frac{2a}{\sigma}\sqrt{2\log b}\right],\qquad b\ge1.
\]
For every $N\ge\kappa_+$, every $m\ge1$ and every $M\ge3$,
\[
 R_\cp^\Theta(N,m)\asymp F_G(N/\kappa_+),\qquad
 R_\aut^\Theta(N,M)\asymp F_G(N/\kappa_+).
\]
Constants depend on $q,a,\sigma,\Gamma$, not on $N,m,M,p_-$. The autonomous upper uses two live states and one terminal prediction. The lower holds for unrestricted history-based stopping and arbitrarily large finite dictionaries. The rare-herald statement is on $N\ge\kappa_+$, not below that feasible comparison range.
\end{theorem}
\begin{proof}
Before the productive Gaussian report the hidden positive probability $r$ lies in $[q,1-q]$. For $y\ge0$, set $t=e^{-2ay/\sigma^2}$ and $c_q=q/(1-q)$, $C_q=1/c_q$. The conditional negative/positive odds are $[(1-r)/r]t\in[c_qt,C_qt]$. Thus for $b_y=\E[E\mid h,Y=y]$,
\[
 \frac{4c_q}{(1+C_q)^2}t\le1-b_y^2\le4C_qt,
 \qquad |b_y-1|^2\le4C_q^2t.
\]
The symmetric calculation applies for $y<0$. Hence the raw task admits the uniform mark
$f(y)=\operatorname{sgn}y$, $w(y)=e^{-2a|y|/\sigma^2}$ by \eqref{eq:ordinary-calibration}. The tie at zero has probability zero. The actual reference acquisition law is the fair mixture
$\nu=\frac12N(a,\sigma^2)+\frac12N(-a,\sigma^2)$. Every full-history productive conditional law lies between $2q\nu$ and $2(1-q)\nu$. This is not a report-row equality: the mixture coefficient $r$ depends on history.

We show uniformly for every $m\ge1$ that $\Qm_m(b)\asymp F_G(b)$. Let $z_b$ be the upper-tail threshold for $|Y|$ with $\nu(|Y|>z_b)=1/(2b)$. Any density $h\le b$ of mass one puts at least half its mass in $|Y|\le z_b$. Discarding the nonnegative codepoint term gives
\[
 \int d_Ch\,d\nu\ge\tfrac12e^{-2az_b/\sigma^2}.
\]
For $b\ge2$, choose $y_b\ge0$ with $\nu(Y\ge y_b)=1/b$, take $h=b\one_{[y_b,\infty)}$ and the single codepoint $+1$. On this set $f=+1$, so the profile is at most $e^{-2ay_b/\sigma^2}$.

For completeness, integration by parts for the standard Gaussian tail gives, for $x>0$,
\[
 \frac{x}{1+x^2}\frac{e^{-x^2/2}}{\sqrt{2\pi}}
 \le\Pp(G\ge x)\le
 \frac1x\frac{e^{-x^2/2}}{\sqrt{2\pi}}.
\]
The upper follows by replacing $1$ in the integral by $u/x$; integration by parts and $u^{-2}\le x^{-2}$ give the lower. These inequalities applied to the two fixed shifted tails imply
\[
 y_b=a+\sigma\sqrt{2\log b}+O\!\left(\frac{\log\log b}{\sqrt{\log b}}\right),\quad
 z_b=a+\sigma\sqrt{2\log b}+O\!\left(\frac{\log\log b}{\sqrt{\log b}}\right)
\]
as $b\to\infty$. Consequently both exponentials are comparable to $F_G(b)$. The remaining bounded $b$ interval has positive continuous quantiles and is handled with constants independent of $m$; a fixed output at the first productive report gives its upper. Finally apply Theorem~\ref{thm:main}. Fixed multiplicative changes of $b$ only change $F_G$ by bounded factors, and $\kappa_+/\kappa_-\le\Gamma$. Two incompatible physical ports and a fixed $+1$ terminal state implement the tail rule. No posterior, hidden transition parameter, or time counter is stored. The mathematical threshold depends only on the known Gaussian and resource certificates; finite-precision evaluation of it requires its own reported numerical accuracy and budget reserve.
\end{proof}

\section{Finite precision with a stopped resource certificate}
\label{sec:numerics}
For a reference finite live alphabet, write $Q$ for its substochastic continuation matrix, $c_i\ge c_*>0$ for reference per-call costs, and $b_i\ge0$ for the expected terminal loss contributed on one call from label $i$. The terminal values themselves are not live states. A proof certificate consists of nonnegative finite vectors
\begin{equation}\label{eq:potentials}
 W\ge c+QW,\qquad U\ge b+QU.
\end{equation}
The implementation may be history dependent and have different one-call quantities $\widetilde Q_h,\widetilde b_h,\widetilde c_h$, but must stay inside the entire certified live alphabet and choose physically legal commands exactly.

\begin{theorem}[Killed-potential implementation bound]\label{thm:perturbation}
Suppose, on every surviving implementation history at live label $i$,
\begin{align}
 ((\widetilde Q_h-Q)W)_i&\le\rho c_i,&0\le\rho&<1,\label{eq:drift-W}\\
 \widetilde b_h(i)+(\widetilde Q_hU)_i&\le U_i+\eta c_i,&\eta&\ge0,\label{eq:drift-U}\\
 \widetilde c_h(i)&\le(1+\zeta)c_i,&\zeta&\ge0.\label{eq:drift-cost}
\end{align}
The implemented observer stops almost surely and satisfies
\begin{equation}\label{eq:implementation}
 \E\sum_{t\le\widetilde\tau}\widetilde c_t
   \le\frac{1+\zeta}{1-\rho}W_{i_0},\qquad
 \E[\text{terminal loss}]
   \le U_{i_0}+\frac{\eta}{1-\rho}W_{i_0}.
\end{equation}
These are uniform bounds if the assumptions are uniform over the unknown model and full raw history.
\end{theorem}
\begin{proof}
The first certificate and \eqref{eq:drift-W} imply
\[
 W_i-(\widetilde Q_hW)_i\ge(1-\rho)c_i.
\]
Set the potential to zero after stopping. Summing to $\widetilde\tau\wedge n$ gives
\[
 (1-\rho)\E\sum_{t\le\widetilde\tau\wedge n}c_{S_{t-1}}
 +\E[W_{S_n}\one_{\{\widetilde\tau>n\}}]\le W_{i_0}.
\]
Keep the nonnegative remaining term until this inequality has been proved. Monotone convergence yields a finite expected reference cost. Its positive minimum implies finite expected duration and almost-sure stopping. Summing \eqref{eq:drift-U} similarly gives
\[
 \E[\text{loss contributed by time }n]
 +\E[U_{S_n}\one_{\{\widetilde\tau>n\}}]
 \le U_{i_0}+\eta\E\sum_{t\le\widetilde\tau\wedge n}c_{S_{t-1}}.
\]
Drop only the nonnegative remaining $U$ term and let $n\to\infty$. Nonnegative loss and almost-sure stopping give the second bound. The actual-cost bound follows separately from \eqref{eq:drift-cost}. Terminal-law total variation alone was not used to bound expected duration.
\end{proof}

\begin{proposition}[Finite comparison implementation of the regular upper]\label{prop:finite-cells}
Suppose the packing data used in Proposition~\ref{prop:geometry} are supplied by finite certified centers, and the current marked coordinates $(f,\sqrt w)$ and those centers can be evaluated to error at most $t/8$ at a declared cost. Assume centers can be approximated within $H$ with that certificate. The regular upper in Corollary~\ref{cor:regular} then has a finite-comparison implementation with the same $B+m$ label bound and the same risk order. For coordinates in a fixed rational box of dimension $k$, precision $p=O(1+\log(1/t))$ suffices; a direct implementation uses $O(mkp)$ program bits, $O(kp+\log m)$ workspace bits, and $O(mkp^2)$ elementary bit operations per productive call, besides the certified coordinate-evaluation cost.
\end{proposition}
\begin{proof}
Use the $j\le m$ disjoint inner rectangles of radius and reliability threshold $t$ constructed in Proposition~\ref{prop:geometry}. Their union has mass at least $1/b$. Replace the locations by certified approximations and use an outward threshold between $3t/2$ and $2t$ for the approximate location distance and reliability, chosen dyadically with its arithmetic error included. Every true inner rectangle is accepted, and every accepted true point has location distance and reliability at most $3t$ from its reported center. Accept the union deterministically, outputting the first matching center. Conditional acceptance probability is at least $\ell/b$, so the same predictable-attempt proof gives cost at most $\kappa_+b/\ell$. The marked distortion of every accepted point is at most $18t^2$, up to the fixed center-approximation factor. No normalization by an unknown exact rectangle mass and no arbitrary density-evaluation oracle is needed.

Store the finitely many rational centers and thresholds in a read-only table. A sequential scan keeps the current coordinates, temporary squared distances and a cell index in workspace. Schoolbook rational comparisons give the displayed bit bounds; all live routing states remain those of $\mathcal G$. The finite set of output values is the center list. At bounded effective budget, accept the first productive report and use one fixed center. This proves the complete range after changing a fixed risk constant. The proposition is conditional on computable certified geometric and report data; it neither computes an arbitrary Borel kernel nor gives an efficient global method for discovering those certificates.
\end{proof}

\subsection*{A complete dyadic certificate}
Take Theorem~\ref{thm:classical} with $K=[0,1]$, independent uniform reference coordinates $(Y,T)$ under $\nu_0$, $\beta=1/4$, $\phi(y,t)=2y-1$, $q=1/4$, $\gamma=1/4$, and gate bounds $p_-=1/8$, $p_+=1/4$. The actual primary laws are not independent of the hidden state; only the primitive reference measure is a product. Here
\[
 c=3/4,\quad C=5/4,\quad a_0=1/20,\quad
 a_v=19/100,\quad k_-=19/300,\quad k_+=3.
\]
Use a three-state observer: gate, acquisition, and terminal value $1/2$. Its ideal acquisition acceptance set is
\[
 A=\{3/8\le Y\le5/8,\ 0\le T\le1/8\}.
\]
The area is $1/32$, so every productive conditional acceptance probability is at least $s_0=3/128$. The reference continuation matrix takes gate success $1/8$ and acquisition acceptance $s_0$. With reference costs $c_G=c_A=1$, its exact cost potentials are
\begin{equation}\label{eq:dyadic-potentials}
 W_G=384,\quad W_A=376,
 \qquad U_G=U_A=1/16.
\end{equation}
Indeed $W_G=1+(7/8)W_G+(1/8)W_A$ and $W_A=1+(1-s_0)W_G$. Every actual ideal row has at least these success probabilities, so satisfies the same super-solutions. On acceptance, $|Y-1/2|+T\le1/4$ and the actual squared audit loss is at most $1/16$, proving the $U$ certificate without knowing the model.

Now let $p=14$, $e=2^{-p}$, and assume the report digitizer supplies $p$-fractional-bit coordinates clipped to $[0,1]$ with deterministic errors at most $e$. Accept only if
\begin{equation}\label{eq:inward-test}
 3/8+e\le\widehat Y\le5/8-e,\qquad \widehat T\le1/8-e.
\end{equation}
Every accepted true report belongs to $A$. Every true report in
\[
 A_{\rm in}=[3/8+2e,5/8-2e]\times[0,1/8-2e]
\]
is accepted. In the coupling using the \emph{same physical report}, a changed routing decision can therefore occur only in $A\setminus A_{\rm in}$. Its conditional probability is at most
\[
 d_e=\frac54\left\{\frac1{32}-(1/4-4e)(1/8-2e)\right\}
      =\frac{10235}{134217728}.
\]
This is a boundary-mass estimate, not a claim that a continuous law and its digitized atomic law are close in total variation. Gate routing is unchanged. Reducing acceptance can increase the acquisition continuation potential by at most $384d_e$, so one valid complete certificate is
\begin{equation}\label{eq:dyadic-errors}
 \rho=\frac{30705}{1048576},\qquad \eta=0,\qquad\zeta=0.
\end{equation}
The $U$ inequality still holds because every accepted true loss is at most $1/16$ and both live $U$ values equal that bound. Thus the implemented observer obeys
\begin{equation}\label{eq:dyadic-final}
 \E\tau\le\frac{384\cdot1048576}{1017871}<400,
 \qquad R\le1/16.
\end{equation}
The strictly positive slack from the physical budget $400$ is
$400-402653184/1017871=4495216/1017871$. A direct conditional-hazard calculation gives the sharper count bound
$9/[\frac34(1/4-4e)(1/8-2e)]$; the looser displayed interval already follows from the killed-potential certificate.

For an all-strategy lower at the same raw budget and alphabet, the two physical ports require two live labels, leaving one output. The effort lower gives at most $400/5=80$ productive calls. Their primary densities are at most $5/4$ relative to the uniform reference, so the terminal law has density at most $100$. For any output $r$, the sublevel
$T^2+(Y-r)^2\le z$ in $[0,1]^2$ has area at most $\pi z/2$. Integrating $\Pp(T^2+(Y-r)^2>z)\ge(1-50\pi z)_+$ yields mean marked distortion at least $1/(100\pi)$. Therefore
\begin{equation}\label{eq:dyadic-risk-interval}
 \frac{19}{30000\pi}\le R_\aut^\Theta(400,3)\le\frac1{16}.
\end{equation}
The lower allows all randomized legal three-state procedures, not only threshold rules. These certified constants are not claimed optimal.

The thresholds in \eqref{eq:inward-test}, in integer units $2^{-14}$, are $6145,10239,2047$. Three 15-bit constants occupy 45 read-only program bits, in addition to a fixed instruction interpreter and the exactly represented output $1/2$. Two 15-bit input registers and a fixed number of control bits suffice as temporary workspace; comparisons take $O(p)$ bit operations per productive report. There are three persistent labels, encoded by two bits, and no stored time counter or retained random seed. If a physical call plus digitization costs at most $t_{\rm raw}(p)$ time and the comparisons cost at most $C_{\rm cmp}p$, the expected execution time is bounded by $400[t_{\rm raw}(14)+14C_{\rm cmp}]$. The promised digitizer accuracy and its cost are part of this implementation contract, not supplied at infinite precision for free.

\section{Causal simulation and a simultaneous risk law}
\label{sec:composition}
We use $\TV(P,Q)=\sup_A|P(A)-Q(A)|$. A causal experiment morphism is a parameter-independent simulator with a strategy lift, a report transformation, a compatible physical-command map, an explicit internal state and clock, and a terminal task readout. The simulator must complete each virtual report before it receives later physical reports. Its error certificate compares the \emph{entire stopped scored transcript}, including the audit and stopping decision, under each lifted legal strategy. This is stronger than comparing a fixed-time terminal marginal.

\begin{proposition}[Typed error and resource composition]\label{prop:composition}
Suppose a simulator from a physical experiment to a target experiment has at most $S$ internal states, uniform stopped-transcript defect at most $\eps$, and separately certified expected physical costs for every lifted procedure. A target observer using at most $M$ states then has a physical implementation with at most $SM$ paired states. For the same nonnegative loss bounded by $L$, its risk is at most the target risk plus $L\eps$. Two such simulators have additive defects and multiplicative state bounds, provided the second certificate applies uniformly to the lifted strategy class and the first simulator's complete interface.

If an additional deterministic output error is at most $\delta$ and the exact prediction-to-audit distance is at most $D$, the pointwise squared-loss change is at most $2D\delta+\delta^2$. Perturbations of stopping or costs instead require a certificate such as Theorem~\ref{thm:perturbation}. No matching lower follows merely from allowing these errors.
\end{proposition}
\begin{proof}
Store the simulator state and observer state as a pair. The command, report and task compatibility requirements make the composition executable by induction over completed virtual calls. No report can be obtained from future data. For a loss in $[0,L]$, integration of its level sets bounds the change in expectation by $L\TV$. Triangle inequality and contraction of total variation under the common second stochastic map give additive errors. The second map must be common and its strategy class valid; otherwise this argument is unavailable. The product of the finite alphabets counts the internal simulator state, with any retained clock or seed included. Expected physical cost is controlled by its own certificate, not by total variation. Finally expand $\|Z-(r+\Delta)\|^2-\|Z-r\|^2$ and use $\|\Delta\|\le\delta$. A one-way simulator only transfers an upper; a converse needs a reverse compatible simulation or an independent decision lower.
\end{proof}

The next result supplies actual alternatives for simultaneous lower terms. It is a selected same-task product experiment, not a universal assertion that every deficiency or numerical tolerance creates unavoidable risk.

\begin{theorem}[Actual calibration and early erasure]\label{thm:joint}
Take either raw realization of Section~\ref{sec:realizations}, with regular geometry and $p_+/p_-\le\Gamma$. Before the core, pay for one call that prepares a fair sign $E$. It reports that sign except on an independent erasure event of probability $2\eps$, $0\le\eps\le1/2$, when it reports $\mathtt{?}$. This virtual interface must finish before the first core call. The core hidden state and all its primitive randomness are independent of $E$ and of the erasure indicator. After commitment, one paid vector audit reports
\[
 (Z,E,\sigma\delta)\in H\times\{-1,1\}\times\{-\delta,\delta\},
 \qquad \sigma\in\{-1,1\}.
\]
The fixed unknown calibration sign $\sigma$ does not enter any pre-audit kernel and is not observed before this audit. Use the sum of the three squared losses, with all three predictions fixed at commitment. Write $N_c=N-2$. For $N_c\ge\kappa_+$ and $M\ge13$,
\begin{equation}\label{eq:joint-law}
 \mathcal R_\aut^\Theta(N,M,\delta,\eps)
 \asymp F(N_c/\kappa_+,M)+\delta^2+\eps.
\end{equation}
The same conclusion holds for checkpoint dictionary size $M$. Constants depend on the raw and geometric certificates, not on $N,M,\delta,\eps$. The early erasure channel has exact causal deficiency $\eps$ to the un-erased early-sign interface under the declared immediate-completion rule.
\end{theorem}
\begin{proof}
On the erased branch every immediate virtual sign is independent of the fair $E$. Against the target's correct pair (early sign, audited $E$), conditional total variation is at least $1/2$: the equality event has target probability one and simulated probability $1/2$. Its branch mass is $2\eps$. Returning a fresh fair sign on erasure attains total variation $\eps$ for the complete scored transcript. On non-erasure the exact sign is passed through. The core cannot improve this immediate simulator: its reports arrive later and, inductively from the product raw kernels and fresh primitive coins, contain no information about an erased $E$. The worst core policy may depend on the reported early tag, but conditional on that tag the same core certificates hold.

For the risk lower, every procedure's core productive count is bounded by its expected core cost, which is at most $N_c$. Its terminal core-mark law is therefore dominated exactly as in Lemma~\ref{lem:domination}, even after conditioning and then averaging over the early tag. Projecting its at most $M$ terminal predictions to the core carrier leaves at most $M$ values. The calibrated geometric lower gives a constant multiple of $F(N_c/\kappa_+,M)$, uniformly over the model. On erasure the best prediction of $E$ has squared risk one, hence its unconditional contribution is $2\eps$. The two unobserved calibration alternatives have identical pre-audit laws. Averaging their risks yields $\delta^2+\E r_\sigma^2\ge\delta^2$; one of them is at least this average. Since the core lower holds under either alternative, the three nonnegative terms hold simultaneously. No supremum/infimum interchange is involved.

For the upper use one initial label and three copies of the two-live-label core program, one for each tag $+, -,\mathtt{?}$. Each copy uses a core dictionary of size
\[
 m=\left\lfloor\frac{M-1}{3}\right\rfloor-2\ge2,
 \qquad 1+3(2+m)\le M.
\]
Predict the reported sign, or zero on erasure, and predict zero for the calibration coordinate. Apply Corollary~\ref{cor:regular} at core budget $N_c$. Since $m$ is bounded below by a fixed positive multiple of $M$ for $M\ge13$, its geometric rate is comparable to that in \eqref{eq:joint-law}. The other two risks are exactly $2\eps$ and $\delta^2$. The preparation call and vector audit are the two additional physical calls, so the total budget is $N_c+2=N$. The same construction has a checkpoint dictionary of at most $M$ triple-valued predictions. The lower for checkpoint is identical. This establishes the stated same-task law, not just separate sharpness examples for separate tasks.
\end{proof}

\appendix
\section{Predictive realization and terminal submeasures on Borel spaces}
\label{sec:flow}
The main theorem is written directly on raw histories. The following results specify how its histories may be replaced by predictive states without assuming a formal manifold. They also extend the terminal-submeasure preservation step beyond a finite driver. Neither result supplies an infinite predictive coordinate as free controller memory.

Fix a declared class $\mathcal T$ of executable future tests, including their command compatibility, costs and terminal readouts. Assume it has a countable determining family $(T_j)$: agreement of their conditional laws implies agreement for every declared test. Put
\[
 q(h)=(\law(T_j\mid h))_{j\ge1}.
\]
The coordinate spaces of probability laws are standard Borel. Assume the actual image $S=q(\mathsf H)$ is Borel and admits a Borel section, that legal command domains are constant on fibers, and that one-step report/update laws for each legal command are determined by $q(h)$ and jointly Borel in command and state. These are realization hypotheses; arbitrary measurable equivalence relations do not automatically possess them. They can be verified from finite raw filtering instruments, a countable determining measurement class and the indicated report-kernel continuity.

\begin{proposition}[Minimality for the declared executable test class]\label{prop:quotient}
Under these hypotheses, $q$ is a measurable realization of the predictive quotient. Its report and update kernel is jointly Borel and yields the same declared future-test laws. Every other Borel statistic determining those laws and admitting the corresponding Borel factorization carries $q$ as a measurable function; no smaller statistic can separate fewer histories while still determining all declared tests.
\end{proposition}
\begin{proof}
The map $q$ is Borel coordinatewise. Its fibers are exactly declared predictive equivalence by determination. Composing a raw one-step conditional kernel with the Borel section gives a Borel kernel on $S$; fiber constancy makes it independent of the chosen section. Push the continued history through $q$ to obtain the predictive update. Induction over finite executable prefixes proves equality of test laws; bounded stopped tests follow by the stated stopping-class approximation when it is part of $\mathcal T$. For another sufficient statistic, equality of its value implies equality of every coordinate of $q$. Its assumed Borel factorization yields the factor map. This proves minimality for the declared class, not an unrestricted control-equivalence theorem.
\end{proof}

For a known prepared model, the usual hidden-state conditional distribution supplies such an endogenous coordinate whenever the tests determine it; otherwise its predictive quotient is the smaller image. For an unknown fixed model, the family of conditional predictive states indexed by $\theta$ is a mathematical comparison object. A modelwise posterior is not a common executable observer. The main upper program uses only its finite label and current report; the full coordinate is used to verify history-uniform inequalities and incurs no hidden implementation claim.

\begin{theorem}[Finite occupation and exact terminal-measure preservation]\label{thm:borel-flow}
Let $S,A,Y$ be standard Borel, let $\alpha$ be a prepared probability measure on $S$, and let $K(ds',dy\mid s,a)$ be a jointly Borel physical transition/report kernel with Borel legal-command and stopping domains. Suppose a finite nonnegative measure $\eta(ds,da)$ is supported on legal commands, and a terminal measure $\zeta(ds,da,ds',dy)$ satisfies
\begin{equation}\label{eq:borel-balance}
 0\le\zeta\le\eta K,\qquad
 \eta_S=\alpha+(\eta K-\zeta)_{s'},
\end{equation}
with $\zeta$ supported on legal stopping. There is a stationary randomized physical policy and a Borel post-report stopping rule whose actual occupation $\bar\eta$ satisfies $\bar\eta\le\eta$ and whose full terminal submeasure is exactly $\zeta$. The policy terminates almost surely with finite expected call count. All nonnegative additive costs can only decrease. The conclusion is for this fixed kernel, not for a collection of separately chosen modelwise policies claimed to be one common simulator.
\end{theorem}
\begin{proof}
Write $\xi=\eta_S$ and disintegrate $\eta(ds,da)=\xi(ds)\pi(da\mid s)$. Take a Borel version $h=d\zeta/d(\eta K)$ clipped to $[0,1]$ and zero outside the legal stopping domain. Define
\[
 Q(s,D)=\int\pi(da\mid s)\int\one_D(s')(1-h(s,a,s',y))K(ds',dy\mid s,a).
\]
It is substochastic and balance becomes $\xi=\alpha+\xi Q$. On $\xi$-null states, command choices and stopping versions may be chosen arbitrarily; the constructed process does not reach such states. Summing \eqref{eq:borel-balance} over the whole state space gives $\zeta(S\times A\times S\times Y)=1$, so terminal mass one is redundant once finiteness and balance hold.

Monotone iteration yields
\[
 \bar\xi=\sum_{n\ge0}\alpha Q^n\le\xi.
\]
In particular $\bar\xi$ is finite and is the actual pre-call occupation of $\pi$ with acceptance $h$. If $d=\xi-\bar\xi$, continuity from below gives $d=dQ$. Hence
\[
 \int(1-Q\one)\,dd=0.
\]
The integrand is nonnegative. A state carrying residual mass has zero acceptance, $d$-almost everywhere. Let $r=d\bar\xi/d\xi\in[0,1]$. The constructed terminal measure is $r(s)\zeta(ds,da,ds',dy)$, and its difference from $\zeta$ has total mass
\[
 \int(1-r(s))\!\int\pi(da\mid s)K(ds',dy\mid s,a)h(s,a,s',y)\,\xi(ds)
   =\int(1-Q\one)\,dd=0.
\]
Thus the terminal measure is preserved exactly. Equivalently, residual circulation is absent on accepting rows and irrelevant on zero-acceptance rows. Finiteness of $\bar\xi(S)$ gives finite expected duration and almost-sure termination. Finally $\bar\eta=\bar\xi\pi\le\xi\pi=\eta$, so every nonnegative additive cost decreases or stays fixed.
\end{proof}

This is the submeasure-preserving form of the classical occupation/Markovization mechanism, now stated for Borel predictive carriers rather than only finite rows. It is adjacent to the extreme finite occupation results of Piunovskiy--Zhang \cite{PZExtreme} and the absorbing-model analyses \cite{PZContinuity,DPAbsorb,DPChatter,DPMixture}. It does not assert that all absorbing occupation sets are weakly compact, nor that this stationary policy has finite controller memory. If two unknown models have different $K$, applying the theorem separately does not produce a common program. The genuinely common finite-state construction in the main theorem is instead Lemma~\ref{lem:rejection}, certified on every raw history.

\section{Scope and separation of the mathematical layers}
The central sufficient-class theorem no longer requires an exact common finite primary row. Its assumptions are an actual acquired-mass envelope, a conditional task comparison characterized by Theorem~\ref{thm:fiber}, an all-strategy effort drift, and a verified finite-state acquisition program. These assumptions are not necessary for every finite-resource experiment. The fiber criterion is necessary and sufficient only for its specified all-prediction calibration mechanism. The reference acquired law is not a formal reachable support; alternative unknown-model laws are compared to it through actual conditional measure inequalities.

The theorem is load-bearing across the two raw families: their state-dependent primary laws and gate probabilities violate fixed-row closure, while the same effort, mass and calibration argument proves their risk curves. The declared finite physical programs, not their potentially infinite posterior trajectories, supply the implementation. The general Borel occupation theorem clarifies exact attainability of a terminal submeasure but does not reduce arbitrary hidden experiments to finite-state controllers.

Physical-command compatibility is invariant under stateless resource-preserving recodings. It does not compute the minimal automaton for an arbitrary protocol, and it is not preserved by adding a free phase-aware primitive or a free external clock. The two-port application has a matching live-state certificate. The numeral $M$ always counts a total alphabet for autonomous observers and a dictionary for checkpoint observers. Temporary workspace, program length, precision, call count, physical time and simulator state remain separate resources.

The regular marked geometric calculation is inherited mathematics, not a newly claimed quantization exponent. New conclusions concern its all-history predictable transfer, the criterion for its task calibration, unknown/history-dependent heralding, physical presentation control, and the actual atomic transition. The finite-bit example provides an explicit correct certificate and an interval, not an optimal arithmetic/compiler theorem. The same-task product theorem uses actual alternatives; general permitted errors produce only certified upper perturbations. No total-variation estimate is promoted to an LDP equivalence or to control of an uncertified mean stopping time.

A near-necessary acquisition classification for arbitrary unknown kernels, matching curves when the conditional mass envelopes fail, a full joint vanishing-$q$ law, efficient global program synthesis, and sharp lower bounds for program/workspace/physical time are not proved. Nor are unrestricted infinite-horizon or unbounded-operator assertions. A standard natural sensing problem with an independently established external open-problem significance is not claimed solved here. The finite-state program and the raw physical families are specified mathematical experiments, and novelty and significance require independent literature and proof assessment.

Historical companion articles are optional provenance and preserve distinct results on conditional pooling, continuation geometry, Bayesian curvature, finite-time allocation and stopped occupation. They are not unstated premises of this article. The main text contains the definitions, central proof, both raw verifications and the error composition needed for its own claims.

\begin{thebibliography}{99}
\bibitem{GTF46} Q. Qi, \emph{General Theta Foundations I: Acquired Geometry and Causal Resource Transfer}, causal-stopping companion, 2026.
\bibitem{GTF51} Q. Qi, \emph{General Theta Foundations I: Acquired Geometry and Causal Resource Transfer}, causal-acquisition companion, 2026.
\bibitem{GTF52} Q. Qi, \emph{General Theta Foundations I: Acquired Geometry and Causal Resource Transfer}, stopped-experiment companion, 2026.
\bibitem{GTF53} Q. Qi, \emph{General Theta Foundations I: Acquired Geometry and Causal Resource Transfer}, stopped-calibration-transfer companion, 2026.
\bibitem{PZExtreme} A. Piunovskiy and Y. Zhang, \emph{Extreme occupation measures in Markov decision processes with an absorbing state}, SIAM J. Control Optim. \textbf{62} (2024), 65--90. DOI: \href{https://doi.org/10.1137/23M1572398}{10.1137/23M1572398}.
\bibitem{PZContinuity} A. Piunovskiy and Y. Zhang, \emph{On the continuity of the projection mapping from strategic measures to occupation measures in absorbing Markov decision processes}, Appl. Math. Optim. \textbf{89} (2024), article 58. DOI: \href{https://doi.org/10.1007/s00245-024-10124-7}{10.1007/s00245-024-10124-7}.
\bibitem{DPAbsorb} F. Dufour and T. Prieto-Rumeau, \emph{Absorbing Markov decision processes and their occupation measures}, SIAM J. Control Optim. \textbf{63} (2025), 676--698. DOI: \href{https://doi.org/10.1137/24M1661959}{10.1137/24M1661959}.
\bibitem{DPChatter} F. Dufour and T. Prieto-Rumeau, \emph{On the Feinberg--Piunovskiy theorem and its extension to chattering policies}, arXiv:2510.04808v1, 6 October 2025.
\bibitem{DPMixture} F. Dufour and T. Prieto-Rumeau, \emph{Absorbing Markov decision processes: geometric properties and sufficiency of finite mixtures of deterministic policies}, arXiv:2512.17511v1, 19 December 2025.
\bibitem{CCGK} K. Chatterjee, M. Chmel\'ik, R. Gupta and A. Kanodia, \emph{Optimal cost almost-sure reachability in POMDPs}, Artificial Intelligence \textbf{234} (2016), 26--48. DOI: \href{https://doi.org/10.1016/j.artint.2016.01.007}{10.1016/j.artint.2016.01.007}.
\bibitem{ACSS} A. Asadi, K. Chatterjee, R. Saona and A. Shafiee, \emph{Limit-sure reachability for small memory policies in POMDPs is NP-complete}, Proc. Forty-first Conf. Uncertainty in Artificial Intelligence, Proc. Mach. Learn. Res. \textbf{286} (2025), 238--256.
\end{thebibliography}

\end{document}
